\documentclass[11pt,twoside,openright]{book}

\usepackage[T1]{fontenc}
\usepackage[utf8]{inputenc}
\usepackage{amsmath,amssymb,amsthm,mathtools,bm}
\usepackage{mathpazo}
\usepackage{microtype}
\usepackage[inner=1.35in,outer=1.1in,top=1.25in,bottom=1.35in]{geometry}
\usepackage{booktabs}
\usepackage{fancyhdr}
\usepackage[hidelinks]{hyperref}
\usepackage{cleveref}

\setlength{\parindent}{1.4em}
\setlength{\parskip}{0pt}
\linespread{1.07}
\emergencystretch=3em

\pagestyle{fancy}
\fancyhf{}
\fancyhead[LE]{\small\thepage\quad\textsc{A Cinema of Individual Editions}}
\fancyhead[RO]{\small\nouppercase{\leftmark}\quad\thepage}
\renewcommand{\headrulewidth}{0.3pt}
\fancypagestyle{plain}{\fancyhf{}\fancyfoot[C]{\small\thepage}\renewcommand{\headrulewidth}{0pt}}

\theoremstyle{plain}
\newtheorem{theorem}{Theorem}[chapter]
\newtheorem{proposition}[theorem]{Proposition}
\newtheorem{lemma}[theorem]{Lemma}
\newtheorem{corollary}[theorem]{Corollary}
\theoremstyle{definition}
\newtheorem{definition}[theorem]{Definition}
\newtheorem{example}[theorem]{Example}
\theoremstyle{remark}
\newtheorem{remark}[theorem]{Remark}

\newcommand{\Capac}{\mathrm{Cap}}
\newcommand{\Rob}{\mathrm{Rob}}
\newcommand{\Inv}{\mathrm{Inv}}
\DeclareMathOperator{\loss}{loss}

\title{\Huge\bfseries A Cinema of Individual Editions\\[1.2ex]
\large Chapter Three: Histories, Folds, and What a Scene Commits To\\[0.4ex]
\large Working Draft}
\author{Flyxion\\[0.3ex]\small Independent Researcher}
\date{November 2026}

\begin{document}

\frontmatter
\maketitle
\tableofcontents
\mainmatter
\setcounter{chapter}{2}

\chapter{Histories, Folds, and What a Scene Commits To}
\label{ch:histories}

The verdict of Chapter~1 ends in a comparison with the script history. Chapter~7 relied on that comparison to separate an error in the film from an act by a character: a loose bolt that appears in a shot is an error if the story says the bolt was tightened, and it is a depiction if the story says that the mechanic skipped it. The history layer decides, and the history as the earlier volume defines it records only what happened. This chapter extends it in two ways. It annotates histories with dramatic function, so that a departure from competence can be marked as intended. It also states a guarantee that a correction cannot delete the evidence on which a later reinterpretation depends. A last section extends the quotient of histories by future equivalence from what can happen to what can be shown, so that a generator tracks a detail such as a wet plate only when an action, a later shot, or an inference refers to it.

The governing example is an engine scene in which an ordinary repair is later understood as diligence, negligence, concealment, or disclosure. It is small enough to count exactly, and it is built so that the same footage supports several readings until a later shot decides among them.

\section{The Engine Scene}\label{sec:engine}

A mechanic works on an engine with three bolts. For each bolt the mechanic either tightens it or skips it. If any bolt is skipped, the mechanic either does not notice, or notices and then reports honestly, or notices and reports falsely that all is well. These choices are the events of the story, and a \emph{history} is a complete record of them. There are eight patterns of tightening and skipping. The pattern in which all three bolts are tightened gives one history, since there is nothing to notice or report. Each of the seven patterns with a skipped bolt gives three histories, namely not noticed, noticed and honestly reported, and noticed and falsely reported. The story world therefore has $1+7\cdot3=22$ histories.

The histories fall into four \emph{readings}, which are the ways a viewer who sees the whole story would describe the repair.

\begin{table}[ht]
\centering
\begin{tabular}{@{}lcl@{}}
\toprule
Reading & Histories & What happened\\
\midrule
Diligence & 1 & all bolts tightened\\
Negligence & 7 & some bolt skipped, not noticed\\
Disclosure & 7 & some bolt skipped, noticed, honestly reported\\
Concealment & 7 & some bolt skipped, noticed, falsely reported\\
\bottomrule
\end{tabular}
\caption{The twenty-two histories of the engine scene and the four readings that partition them.}\label{tab:readings}
\end{table}

The script history $h^\ast$ is the one the production intends: the first two bolts are tightened, the third skipped, and the mechanic notices and reports falsely. It is a concealment.

\section{Annotating Histories with Dramatic Function}\label{sec:annotate}

The history $h^\ast$ departs from competence. A competent mechanic tightens all three bolts, and the competent history is the first of Table~\ref{tab:readings}. The departure is an act of the character, and it is also something the film has to be allowed to show. If the history layer records only what happened, a film that shows the third bolt loose depicts something that differs from the competent baseline, and a report that this is a continuity error cannot be answered by the record.

An \emph{annotation} assigns each event of a history a dramatic function. The functions used here are \emph{competent}, for an action that follows the baseline, and \emph{lapse}, \emph{noticing} and \emph{concealing} for departures that the story intends. In $h^\ast$ the skipping of the third bolt is a lapse, the notice is a noticing, and the false report is a concealing. The annotation does not change which events happened. It changes which departures from the baseline are errors in the film and which are acts of its characters.

\begin{definition}[Verdict with annotation]\label{def:annverdict}
The showing is first compared with the script history. If the event it displays differs from the event of the script history, it is an \emph{error}, whatever the annotation says, since an intended lapse can still be misdepicted. If the event it displays is the event of the script history and that event departs from the baseline, the annotation decides: the showing is a \emph{depiction of an act} when the annotation marks the departure as one the story intends. A report that a depiction of an act is an error is rejected.
\end{definition}

In the engine scene, a report that the third bolt is shown loose in the close shot is rejected, since the displayed state matches $h^\ast$. The same report is upheld as an error if the script history is the all-tightened history and no annotation exists. The shot is the same in both cases and the history differs.

The cost of getting this wrong is counted in editions. The scene as shot has $126$ class tuples, from two classes at the opening, seven at the bolt close-up, and three each at the glance and the handoff. Of these, $32$ are consistent with $h^\ast$, in the sense that no shot in them contradicts it. If the close shot of the loose third bolt is removed as an error, the number of editions consistent with $h^\ast$ falls from $32$ to $24$, and it falls by editions that were telling the story correctly.

\section{What a Scene Commits To}\label{sec:commits}

Each shot class shows something, and what it shows rules out some histories. A class that cuts to the wide view of the bench rules out nothing. A close-up of the first bolt, tightened, rules out every history in which that bolt was skipped. The \emph{commitment} of a class is the set of histories consistent with what it shows, and the commitment of an edition is the intersection of the commitments of its classes, which is the set of histories that the edition leaves open. This is the fold of the earlier volume taken from histories to what can be shown: the edition's commitment is the fiber of the fold over the edition.

The relevant properties of an edition are properties of its commitment $C(e)$.

\begin{definition}[Properties of an edition]\label{def:props}
An edition $e$ is \emph{consistent} if $C(e)$ is nonempty and \emph{faithful} if $h^\ast\in C(e)$. It is \emph{open} to a reading if $C(e)$ contains a history of that reading. It is \emph{open} if it is open to all four readings. It is \emph{recoverable as} a reading if $C(e)$ is nonempty and lies wholly within that reading, so that a viewer who has seen the edition can say which reading holds.
\end{definition}

For the engine scene $112$ of the $126$ editions are consistent, and $32$ are faithful. A faithful edition can be recoverable only as the reading that $h^\ast$ belongs to, which is concealment. Eight of the $32$ are recoverable as concealment, and six are open to all four readings. No edition is both. An edition that is open to diligence, whose commitment contains the history with every bolt tightened, cannot contain the glance in which the mechanic notices, and an edition that supplies the glance cannot be open to diligence. A scene therefore cannot both leave every reading available and make one of them recoverable. The production's choice, edition by edition, is which of the two to give the viewer, and a production that wants both ambiguity and a later resolution has to place them in different shots, with the evidence in the later one.

\section{A Correction Cannot Delete the Evidence}\label{sec:guarantee}

The engine scene is built so that the reading of concealment depends on two classes: the glance in which the mechanic notices, and the handoff in which the report is made without comment. Every edition that is recoverable as concealment uses both. The eight recoverable editions differ only in which of the four close-up classes consistent with $h^\ast$ they use and in which class of the opening shot, and in neither are they constrained by the evidence.

\begin{proposition}[Preservation of evidence]\label{prop:evidence}
Let $E$ be the set of editions that are recoverable as a reading $R$, and let an amendment remove a set of classes. The editions recoverable as $R$ after the amendment are those of $E$ that use none of the removed classes. In particular, recoverability survives the amendment if and only if some edition of $E$ avoids every removed class, and an amendment that removes a class used by every edition of $E$ leaves none.
\end{proposition}

\begin{proof}
An edition's commitment depends only on its classes, so an edition that avoids the removed classes keeps its commitment and its recoverability, and one that uses a removed class is no longer an edition of the production.
\end{proof}

The numbers follow. Withdrawing the glance in which the mechanic notices leaves none of the eight editions recoverable as concealment, and so does withdrawing the clean handoff. Withdrawing one class of the opening shot leaves four of the eight, since the opening is irrelevant to the evidence. The withdrawal of the glance also halves the number of faithful editions, from $32$ to $16$. The glance can be withdrawn for reasons that have nothing to do with the story, for example because the rights to the shot are lost or the take is damaged, and when it is, the later reinterpretation depends on a class that no longer exists.

The guarantee that the production should give is a condition on amendments. A correction answering a report is \emph{evidence-preserving} for a reading if at least one edition recoverable as that reading avoids every class the correction removes. The test is on the removed set as a whole. It is not enough that no single removed class is used by every recoverable edition: if two recoverable editions use alternative classes $a$ and $b$ at one slot, neither class occurs in every recoverable edition, and removing both destroys recoverability. For an amendment that is the agreement rule with the script history, the condition holds for the reading to which $h^\ast$ belongs whenever some faithful edition is recoverable as it, by the proposition below.

\begin{proposition}[Agreement with the script preserves faithful editions]\label{prop:agree}
An amendment that removes exactly the classes whose commitments exclude $h^\ast$ leaves every faithful edition in place, and therefore preserves recoverability of the reading to which $h^\ast$ belongs.
\end{proposition}

\begin{proof}
A faithful edition has $h^\ast$ in its commitment, so none of its classes excludes $h^\ast$. The amendment removes only classes that exclude $h^\ast$, so every faithful edition survives.
\end{proof}

The classes that contradict $h^\ast$ in the engine scene are five: two classes at the bolt close-up that show the first or second bolt loose, the class that shows the third bolt tight, the glance in which the mechanic does not look, and the honest handoff. Agreement with the script removes those and nothing else. The danger is an amendment adopted on a mistaken verdict, as in the example of the loose third bolt, which removes a class that is consistent with $h^\ast$ and thereby removes editions that carried the story. Chapter~7's verdict, with the annotation of this chapter, is what prevents it.

\section{Noise Overcapture}\label{sec:noise}

A scene has details that matter and details that do not. The engine scene is shot in a workshop with four plates on a rack, each wet or dry, and nothing in the story refers to them. A generator that tracked every detail of the shot would carry each plate's state as part of the history, and the number of histories would grow with the details.

The earlier volume handles this by a quotient of histories by future equivalence: two prefixes of a history are identified if they allow the same futures. This chapter extends it from what can happen to what can be shown. Two partial histories are \emph{equivalent for the future} if every later shot, action, and inference that refers to the history treats them alike. A plate matters in exactly three ways: an action depends on it, such as the mechanic slipping on a wet plate, a later shot shows it, or an inference uses it, such as a viewer deciding that the mechanic was washing up and left a bolt for later. Unless one of these refers to the plate, its state can be dropped from the history, and the generator is free to render it any way.

\begin{table}[ht]
\centering
\begin{tabular}{@{}lr@{}}
\toprule
Plates referred to by the future & Classes of histories\\
\midrule
0 & 22\\
1 & 44\\
2 & 88\\
3 & 176\\
4 (all tracked) & 352\\
\bottomrule
\end{tabular}
\caption{Classes of histories of the engine scene, with four plates each wet or dry, as the number of plates referred to by a later action, shot, or inference grows.}\label{tab:noise}
\end{table}

The count doubles with each plate that a later shot refers to, and reaches $22\cdot2^4=352$ only if all four are tracked. A generator that tracks what the future refers to carries $22$ classes if nothing refers to the plates. It carries $44$ if a later shot shows one of them wet, and the state of that plate then enters every history from the moment the plate's state is fixed. This is the same distinction between what must be remembered and what need not be that Chapter~4 drew with carry width, and here it applies to the history layer.

\section{What the Chapter Gives the Others}\label{sec:gives3}

Chapter~7's verdict uses the annotated history: the third step of the verdict compares the verified appearance with the script history and consults its annotation to decide whether a departure is an error or an act. Chapter~8's footprint and correction event take their evidence requirement from \Cref{prop:evidence}: a correction is formable only if it is evidence-preserving for the readings the production has promised to keep available. Chapter~14, on sequels and reinterpretation, begins from the property that a scene can leave a reading open or make it recoverable, and not both.

The script \texttt{verify\_ch03.py} reproduces each figure of this chapter.

\begin{thebibliography}{99}

\bibitem{flyxion2026gf}
Flyxion.
\newblock \emph{Generative Film: Categories, Sheaves, and Types for Cinema}.
\newblock Unpublished monograph, October 2026.

\end{thebibliography}

\end{document}
